% =============================================================================
\chapter{Introduction}
\label{ch:introduction}
% =============================================================================

\section{Overview}

ShepherdLab LW~FEA-OPT is a lightweight, pure-MATLAB toolkit for
finite element analysis and stress-ratio thickness optimisation of
two-dimensional frame and truss structures.  It is designed to be readable,
hackable, and easy to extend.

The core capabilities are:

\begin{itemize}[nosep]
    \item 2-D Euler--Bernoulli frame elements with axial, shear, and bending
          degrees of freedom.
    \item Linear and nonlinear (Newton--Raphson) large-deflection analysis.
    \item Stress-ratio thickness optimisation for uniform stress distribution.
    \item Truss subdivision for organic topology optimisation.
    \item Geometry import from CSV and IGES files.
    \item Image-processing-based profile export for CAD integration.
\end{itemize}

\section{Installation}

Clone or download the repository, then add the source directories to the
MATLAB path by running the setup script:

\begin{lstlisting}
run('path/to/lw-feaopt/setup.m')
\end{lstlisting}

This adds \texttt{src/}, \texttt{io/}, and \texttt{examples/} to the path and
creates the \texttt{profiles/} output directory if it does not already exist.

\section{Requirements}

\begin{itemize}[nosep]
    \item MATLAB R2019b or later.
    \item \textbf{Image Processing Toolbox} --- required for profile export
          (\texttt{bwboundaries}, \texttt{poly2mask}, \texttt{imclose}).
    \item Curve Fitting Toolbox --- optional; enables \texttt{csaps} smoothing
          splines in the importers and exporter.  If unavailable, the code
          falls back to \texttt{pchip} interpolation.
\end{itemize}

\section{Quick Start}

Two demo scripts are included:

\begin{lstlisting}
demo_beam_optimization      % curved beam from an IGS file
demo_truss_optimization     % Warren truss with element subdivision
\end{lstlisting}

Each script follows the same four-stage workflow:

\begin{enumerate}[nosep]
    \item \textbf{Define geometry} --- import a curve or build a truss
          programmatically, assign material and section properties, and
          specify boundary conditions.
    \item \textbf{Analyse} --- run linear and/or nonlinear FEA.
    \item \textbf{Optimise} --- iterate with the stress-ratio method until
          convergence.
    \item \textbf{Export} --- rasterise the optimised structure and export
          spline-ready boundary loops for SolidWorks.
\end{enumerate}

\section{Project Layout}

\Cref{tab:layout} summarises the directory structure.

\begin{table}[htbp]
\centering
\caption{Repository directory structure.}
\label{tab:layout}
\begin{tabularx}{\textwidth}{lX}
\toprule
\textbf{Path} & \textbf{Contents} \\
\midrule
\texttt{src/}       & Core classes: \texttt{StructureGeometry},
                      \texttt{ElementMatrices}, \texttt{AnalysisState},
                      \texttt{FEAAnalysis}, \texttt{LoadCase},
                      \texttt{StructurePlotter}, \texttt{optimizeThickness},
                      \texttt{subdivideTruss}. \\
\texttt{io/}        & Geometry importers (\texttt{importCSVCurve},
                      \texttt{importIGSCurve}) and the profile exporter
                      (\texttt{exportProfile}). \\
\texttt{examples/}  & Runnable demo scripts. \\
\texttt{profiles/}  & Auto-created output directory for exported CAD profiles. \\
\texttt{docs/}      & This documentation ({\LaTeX} source). \\
\bottomrule
\end{tabularx}
\end{table}

\section{Conventions}

Throughout this document:

\begin{itemize}[nosep]
    \item Vectors are set in bold italic: $\vect{u}$, $\vect{f}$.
    \item Matrices are set in bold upright: $\mat{K}$, $\mat{T}$.
    \item Transpose is denoted $(\cdot)\T$.
    \item Scalar quantities use standard italic: $E$, $A$, $I$, $L$.
    \item SI units are used throughout unless otherwise stated.
    \item ``Element'' and ``member'' are used interchangeably.
    \item Node DOFs are ordered $[u_x,\; u_y,\; \theta]$ at each node.
\end{itemize}